% =============================================================
%  数学家 LaTeX 文档共用导言区
%  被 tex/<Name>/<Name>.tex 引入
% =============================================================

\documentclass[11pt,a4paper]{article}

% ---- 编码与字体（xelatex）----
\usepackage{fontspec}
\defaultfontfeatures{Ligatures=TeX}
% 使用系统自带字体（macOS：Times New Roman / Helvetica / Menlo；Linux 可用相似替代）
% 这几个名字在绝大多数 macOS 上都存在；若缺失可在 common/preamble.tex 改为你喜欢的字体
\IfFontExistsTF{Times New Roman}{\setmainfont{Times New Roman}}{}
\IfFontExistsTF{Helvetica}{\setsansfont{Helvetica}}{}
\IfFontExistsTF{Menlo}{\setmonofont{Menlo}}{}
% 不做 CJK 专用字体设定；如文本含中文，可自行加载 xeCJK
\usepackage{newunicodechar}

% ---- 排版与版面 ----
\usepackage[margin=2.3cm]{geometry}
\usepackage{microtype}
\usepackage{setspace}
\onehalfspacing
\setlength{\parskip}{0.5em}
\setlength{\parindent}{0pt}

% ---- 数学 ----
\usepackage{amsmath,amssymb,amsthm}
\usepackage{mathtools}

% ---- 表格 ----
\usepackage{booktabs}
\usepackage{longtable}
\usepackage{array}
\usepackage{tabularx}
\usepackage{multirow}

% ---- 图像 ----
\usepackage{graphicx}
\usepackage{float}
\graphicspath{{images/}}  % 本地图片默认路径
% 新版 LaTeX (>=2019-10) 的 graphicx 已原生支持包含点号/下划线等文件名，
% 不再需要 grffile；引入 grffile 反而会引起 \iffalse 错误。

% ---- 超链接 / 引用 ----
\usepackage[hidelinks,breaklinks=true]{hyperref}
\hypersetup{
  colorlinks=true,
  linkcolor=black,
  urlcolor=blue!60!black,
  citecolor=black,
}
\usepackage{url}
\usepackage{xurl}  % 允许 URL 在任意位置断行，避免 \href@split 超长错误

% ---- 排序/列表 ----
\usepackage{enumitem}
\setlist{topsep=0.2em,itemsep=0.1em,parsep=0.1em}

% ---- pandoc 兼容所需的若干命令 ----
% pandoc --to=latex 有时会用到 \tightlist / \passthrough / \pandocbounded 等
\providecommand{\tightlist}{%
  \setlength{\itemsep}{0pt}\setlength{\parskip}{0pt}}
\providecommand{\passthrough}[1]{#1}
% pandoc 3.x 使用 \pandocbounded 包裹 includegraphics，以防超宽
\providecommand{\pandocbounded}[1]{%
  \makebox[\linewidth][c]{\resizebox{!}{!}{#1}}}

% ---- 段落盒子 ----
\usepackage{framed}
\usepackage{xcolor}

% ---- 页眉页脚 ----
\usepackage{fancyhdr}
\setlength{\headheight}{15pt}
\addtolength{\topmargin}{-3pt}
\pagestyle{fancy}
\fancyhf{}
\fancyhead[L]{\nouppercase{\leftmark}}
\fancyhead[R]{\thepage}
\renewcommand{\headrulewidth}{0.3pt}

% ---- 标题 ----
\usepackage{titling}

% ---- pandoc 可能用到的语言命令兜底（避免加载庞大的 polyglossia/babel）----
% 注意：hyperref 会偷偷定义一个要求加载 babel 的 \foreignlanguage，
% 这里强制覆盖成简单的"原样输出第二参数"。
\AtBeginDocument{%
  \providecommand{\foreignlanguage}[2]{#2}%
  \renewcommand{\foreignlanguage}[2]{#2}%
  \providecommand{\textlang}[2]{#2}%
  \providecommand{\lang}[1]{}%
}

% ---- 元数据命令（由 metadata.tex 填充）----
\providecommand{\MathematicianName}{Unknown}
\providecommand{\MathematicianBirth}{?}
\providecommand{\MathematicianDeath}{?}
\providecommand{\MathematicianNationality}{}
\providecommand{\MathematicianFields}{}
\providecommand{\MathematicianAwards}{}
\providecommand{\MathematicianDescription}{}
\providecommand{\MathematicianWikipedia}{}

% 信息框
\newcommand{\InfoBox}{%
  \begin{center}
  \begin{tabular}{@{}p{0.28\textwidth}p{0.62\textwidth}@{}}
    \toprule
    \textbf{Born}          & \MathematicianBirth \\
    \textbf{Died}          & \MathematicianDeath \\
    \textbf{Nationality}   & \MathematicianNationality \\
    \textbf{Fields}        & \MathematicianFields \\
    \textbf{Awards}        & \MathematicianAwards \\
    \textbf{Wikipedia}     & \url{\MathematicianWikipedia} \\
    \bottomrule
  \end{tabular}
  \end{center}
}
